\documentclass[12pt]{article}
\usepackage[utf8]{inputenc}
\usepackage[T1]{fontenc}
\usepackage{hyperref}
\title{Kurrama Language Resources for K-12 Curriculum Development}
\author{Compiled by Claude Sonnet 5 (claude-sonnet-5) for Andr\'as Kornai}
\date{}
\begin{document}
\maketitle

\section{Introduction}
Kurrama (ISO 639-3 vku, Glottocode kurr1243) is a Pama-Nyungan language of the Pilbara, Western Australia, with roughly ten speakers recorded in 2005 and classed as critically endangered. It is often described as close to Yindjibarndi, though speakers regard it as distinct.

\section{Primary Educational Resources}
No primers, textbooks or school programmes specific to Kurrama were found.

\section{Linguistic and Reference Materials}
\begin{itemize}
\item Peter M. Hill (2011), \emph{Morphology and Sentence Construction in Kurrama}, University of Western Australia dissertation (reference record): \url{https://glottolog.org/resource/reference/id/320244}
\item Wangka Maya Pilbara Aboriginal Language Centre page on Kurrama, which notes a 2006 dictionary and sketch grammar: \url{https://www.wangkamaya.org.au/pilbara-languages/kurrama}
\item Wangka Maya resources and dictionaries page (dictionaries are copyrighted; permission is needed to reproduce): \url{https://www.wangkamaya.org.au/about-languages/resources-and-dictionaries}
\item Community-built Kurrama--English word list on Glosbe (not a scholarly source): \url{https://glosbe.com/vku/en}
\item Overview: \url{https://en.wikipedia.org/wiki/Kurrama_language}
\end{itemize}

\section{Religious and Community Resources}
Wangka Maya is the regional language centre. No religious materials were found.

\section{Digital and Audio Resources}
No audio or video collections were found in search results.

\section{Curriculum Development Context}
The language is not known to be taught in schools. Any programme would depend on Wangka Maya, community permission and the existing dictionary and grammar, whose reuse requires written consent.
\end{document}